\documentclass{article}
\usepackage[utf8]{inputenc}
\usepackage[spanish]{babel}
\usepackage{amsmath}
\usepackage{circuitikz}
\usepackage{float}

% Solución para el conflicto entre circuitikz y babel en español
\usetikzlibrary{babel}


\title{Prácticas de electricidad}
\date{}

\begin{document}

\maketitle

\begin{figure}[H]
    \centering
    \begin{circuitikz}[american] % Eliminamos las opciones de escalado temporalmente para depurar
        % Main path: Battery -> R1 -> Parallel Block 1 (R2, R3) -> R4 -> Parallel Block 2 (R6, R7) -> R5 -> Battery
        \draw (0,0) to[battery, l=$V$, invert] (0,3) % Batería
            to[R, l_=$R_1$, v^=$V_1$] (2,3) node[right] {A} % R1 y nodo A
            coordinate (A_end); % Coordenada para el final de R1 / inicio del bloque paralelo 1

        % Bloque paralelo 1 (R2, R3) - Horizontal
        \coordinate (B_start) at (4,3); % Coordenada para el final del bloque paralelo 1 / inicio de R4
        \draw (A_end) -- (A_end |- 0,4) to[R, l_=$R_2$] (B_start |- 0,4) -- (B_start); % Rama superior (R2)
        \draw (A_end) -- (A_end |- 0,2) to[R, l_=$R_3$] (B_start |- 0,2) -- (B_start); % Rama inferior (R3)
        \node[left] at (B_start) {B}; % Nodo B

        % R4
        \draw (B_start)
            to[R, l_=$R_4$, v^=$V_3$] (6,3) node[right] {C} % R4 y nodo C
            coordinate (C_end); % Coordenada para el final de R4 / inicio del bloque paralelo 2

        % Bloque paralelo 2 (R6, R7) - Horizontal
        \coordinate (D_start) at (8,3); % Coordenada para el final del bloque paralelo 2 / inicio de R5
        \draw (C_end) -- (C_end |- 0,4) to[R, l_=$R_6$] (D_start |- 0,4) -- (D_start); % Rama superior (R6)
        \draw (C_end) -- (C_end |- 0,2) to[R, l_=$R_7$] (D_start |- 0,2) -- (D_start); % Rama inferior (R7)
        \node[left] at (D_start) {D}; % Nodo D

        % R5 y retorno a la batería
        \draw (D_start)
            to[R, l_=$R_5$, v^=$V_5$] (10,3) % R5
            -- (10,0) to[short, i_>=$I_T$] (0,0); % Cierre del circuito

        % Etiquetas de voltaje V2, V4
        \draw[<->] (A_end |- 0,5) -- (B_start |- 0,5) node[midway, above] {$V_2$};
        \draw[<->] (C_end |- 0,5) -- (D_start |- 0,5) node[midway, above] {$V_4$};
        
    \end{circuitikz}
    \caption{Aplicación de ley de tensión de Kirchhoff a un circuito serie-paralelo.}
    \label{fig:19-2}
\end{figure}

un circuito cerrado. También se aplica la ley en el caso de que el circuito sea serie-paralelo (fig. 19-2). Aquí, \( V = V_1 + V_2 + V_3 + V_4 + V_5 \), donde \( V_1, V_3, V_5 \) son las caídas de tensión en \( R_1, R_4, R_5 \), respectivamente. \( V_2, V_4 \) son las tensiones de los circuitos en paralelo entre los terminales \( A, B, y \) entre \( C, D, \) respectivamente.

Expresada en palabras, la ecuación (19-5) enuncia que en un circuito o bucle cerrado, la tensión aplicada es igual a la suma de las caídas de tensión en el circuito.

El concepto de signos algébricos o polaridad es útil para resolver los problemas de redes eléctricas. Un ejemplo aclarará el convenio empleado en la asignación de un signo + o — a una tensión en un circuito.

Consideremos el circuito de la figura 19-3 y que realmente la corriente es un flujo de electrones que se mueven desde un potencial negativo a uno positivo. La flecha de la figura 19-3 indica el sentido de la corriente, y los signos — y + indican lo siguiente: el punto \( A \) es negativo con respecto al \( B \); el punto \( B \) es negativo con respecto al \( C \), etc. Esto está de acuerdo con nuestro concepto de corriente en este circuito.

Ahora para establecer el signo algébrico de las tensiones en el circuito cerrado nos movemos en el sentido admitido de la corriente*. Consideremos como positiva cualquier fuente de tensión o caída de tensión cuyo terminal + (positivo) sea el primero alcanzado, y como negativa cualquier fuente de tensión o caída cuyo terminal — (negativo) sea el primero alcanzado. Partiendo del punto \( A \) en la figura 19-3 y moviéndonos en el sentido de la corriente tenemos — \( V_1, V_2, V_3, V_4, V_5 \) . Teniendo siempre presente este convenio, podemos generalizar la ley de tensión de Kirchhoff, como sigue:

La suma algébrica de las tensiones en un circuito cerrado es igual a cero.

Aplicando el convenio de los signos y la ley de Kirchhoff al circuito cerrado de la figura 19-3 y partiendo del punto \( A \), podemos describir lo siguiente:

\[
V_1 - V_2 - V_3 - V_4 + V = 0 \quad (19-6)
\]

¿Es consistente esta ecuación con la ec. (19-5)? Sí, porque pasando los términos del segundo miembro de la ecuación (19-5) al primer miembro, tendremos \( V - V_1 - V_2 - V_3 - V_4 = 0 \), cuyo resultado es idéntico al de la ecuación (19-6).

\begin{figure}[H]
    \centering
    \begin{circuitikz}[american]
        \draw (0,0) node[below] {E}
            to[battery, l^=$V$, invert] (0,3) node[left] {F}
            to[R, l_=$R_1$, v^>=$V_1$] (2,3) node[above] {A}
            to[R, l_=$R_2$, v^>=$V_2$] (5,3) node[above] {B}
            to[R, l_=$R_3$, v^>=$V_3$] (8,3) node[above] {C}
            to[R, l_=$R_4$, v^>=$V_4$] (11,3) node[above] {D}
            to[short] (11,0)
            to[short, i_>={corriente}] (0,0);
    \end{circuitikz}
    \caption{Asignación convencional de polaridad a las tensiones en un circuito cerrado.}
    \label{fig:19-3}
\end{figure}

* Este sentido es el del flujo de electrones, contrario al convencional de corriente eléctrica que le fue arbitrariamente atribuida antes de descubrirle su naturaleza electrónica (N. del T.).

\end{document}